\section{Related Work}
\label{sec:related_work}

\noindent\textbf{Knowledge and reasoning in language models}. Numerous work finds that transformer language models, even SoTA ones such as GPT-4, struggle in implicit reasoning over their parametric knowledge~\citep{talmor2020olmpics, kassner2020pretrained, rogers-etal-2020-primer, rae2021scaling, press-etal-2023-measuring, allenzhu2023physics, yang2024large}, suggesting their limitations in inducing structured and compressed representations of facts and rules during training. A series of efforts try to understand transformer's knowledge and reasoning through controlled experiments~\citep{kassner-etal-2020-pretrained, prystawski2023why, dziri2023faith, wang2024understanding}, which is also our focus. We find that transformers can learn implicit reasoning over knowledge through grokking, and characterize the connection between the acquired systematicity level and the inductive bias of transformer.

\noindent\textbf{``Chain-of-Thought'' and verbalized reasoning}. A series of studies prompt/fine-tune language models to verbalize (i.e., generate) the intermediate knowledge and reasoning steps~\citep{wei2022chain,wang-etal-2022-iteratively,zelikman2022star,sun2023recitationaugmented,liu-etal-2023-crystal, zelikman2024quietstar} during inference, which has been shown to improve performance especially for large models with strong generation capabilities. There are also theoretical results showing the benefits of such verbalizations~\citep{feng2023towards, li2024chain}. Our focus here on implicit reasoning is orthogonal, and it is an interesting open problem to have principled understandings of the role of such verbalizations in reasoning problems, and also develop methods that can decide the appropriate balance between implicit and explicit reasoning to handle challenging problems with large intrinsic complexity. Relatedly, recent work also finds that explicit verbalizations could be a useful medium for teaching models to reason implicitly via distillation or curriculum~\citep{deng2023implicit, deng2024explicit}. 

\noindent\textbf{Grokking} is first discovered by Power et al.~\citep{power2022grokking} on a set of small algorithmic reasoning tasks. The intriguing phenomenon inspired follow-up works proposing different explanations and expanding the set of tasks where grokking is observed~\citep{thilak2022slingshot, liu2022towards, davies2022unifying, notsawo2023predicting, nanda2023progress, varma2023explaining, merrill2023a, murty-etal-2023-grokking, liu2023omnigrok, zhu2024critical, huang2024unified}. To our knowledge, we are the first work to observe grokking in the domain of knowledge-based reasoning, and our controlled experiments suggest potential corrections of prior hypotheses based on critical data size. Our formulation of rule induction from atomic and inferred facts is general, which we hope could inspire future work on understanding grokking and generalization in deep learning.

\noindent\textbf{Analyzing the inner workings of neural models}. Recent work tries to open up the ``black box'' of neural models through a wide range of techniques; see survey in Ferrando et al.~\citep{ferrando2024primer}. We apply causal tracing~\citep{vig2020investigating, meng2022locating, hanna2023how, wang2023interpretability, feng2024how} and logit lens~\citep{nostalgebraist2020, geva-etal-2022-transformer, yang2024large} to discover interpretable circuits in the model to understand the grokking process and how/why generalization happens.

\noindent\textbf{Parametric and non-parametric memory}. Our focus in this work is on parametric memory in language models, and an orthogonal direction is to enhance models with non-parametric memory, such as extending the effective context length~\citep{xiong2023effective, openai_devday_2023, fu2024data, google2024gemini} and augmenting with retrieval~\citep{guu2020retrieval, lewis2020retrieval, borgeaud2022improving, tay2022transformer, zhong-etal-2022-training, min-etal-2023-nonparametric}. The two types of memory are largely complementary to each other---parametric memory has its unique ability to compress and integrate information but is also inevitably lossy and subject to hallucination, while non-parametric memory is lossless and could also provide attribution. Similarly for humans---a human acquires expertise in a domain by acquiring and structuring knowledge in the brain (parametric), but he/she also wouldn't memorize all pieces of details and could refer to the source when necessary (non-parametric). How to decide the tradeoff between parametric and non-parametric memory (or, how to define the objective for such a tradeoff) is another interesting open problem for future work.
